\silentchapter{To the Reader}{to the reader}
\begin{widebody}
\begin{foreordbody}

The book has seven parts.

The \textit{Glossary} defines the technical terms in alphabetical order. It is for consulting along the way, not for reading first.

The \textit{Tractatus} is axiomatic. It builds a conceptual apparatus from the ground up, through numbered propositions in eight sections. The last section is the generative model, where grammar, landscape, and light cone are united in one equation. The propositions are tagged by type: definition (\textit{d}), theorem (\textit{t}), postulate (\textit{a}), observation (\textit{o}), illustration (\textit{i}).

The \textit{Foundation} shows that the framework is not isolated. The same formalism is empirically demonstrated across fourteen substrate levels, from gene regulatory networks to civilisations. The research programme that follows is sketched in four points.

The \textit{Afterword} is material and programmatic in one flow. From the caterpillar and the disposable cup, via Palantir and the attention economy, to the practical demands on the field.

\textit{Three letters} address the designer, the student, and the academy directly.

The \textit{References} are numbered. The superscript numerals in the propositions and the afterword point here.

The \textit{Notes} comment on individual propositions, place them in the tradition, and list the falsification conditions.

The academic will find most of what matters in the tractatus, the foundation, and the notes. The student will find herself in the afterword. The designer in practice should read the last page first, and begin again.

\end{foreordbody}
\end{widebody}
